\section{Discussion}
\label{sec:discussion}

\emph{What was shown.} Proved: the forward process has the Fokker--Planck equation~\eqref{eq:fp} for any flow; a stochastic equation with the physical drift needs $r \geq \kappa$ at every $\sigma$ (Proposition~\ref{prop:sde}); and in the divergence form errors of learned fields enter the critical rate only linearly, so that errors orthogonal to the true fields leave it exact while the direct form is biased by the square of the score error (Proposition~\ref{prop:linear}). Observed: small networks trained by denoising score matching read the margin $(\kappa - R)/\sigma$ to within $2.2\%$ in two dimensions and $1.4\%$ in four; in four dimensions they agree with the formula for the gap, including its finite-resolution correction, to within $2\%$; in the band $R < r < \kappa$ the entropy grows as predicted while the stochastic equation with the physical drift fails; and the apparent knee below $\kappa$ is predicted by the shape of the missing diffusion.

\emph{Two arrows.} The two thresholds answer different questions. Threshold (II) asks whether the observer loses information: it depends on the state and on $\sigma$, and its margin is the resolution gap. Threshold (I) asks whether the loss can be produced by noise added to the physical dynamics: it depends only on the flow. The band between them shows that the coarse-grained arrow of time does not need a noisy dynamics behind it. For phase-space diffusion models, threshold (I) has a practical reading: a model that uses the physical flow as its drift must choose a noise rate of at least $\kappa$, the largest strain rate, whatever the noise level. Below it, the probability-flow equation remains available, because it requires no positive diffusion.

\emph{The estimator matters more than the network.} The same networks miss the margin by $14\%$ with the plain form and by $1\%$ with the divergence form. Any quantity that is a small difference of two large expectations of learned fields, as the margin is, should be computed in a form in which the learned fields appear once and are tested against the true law. The integration-by-parts identity of Proposition~\ref{prop:linear} achieves this for the entropy rate; it applies to any score-based model whose samples come from the true forward process.

\emph{Limitations.} The networks are small, written in NumPy and trained on one CPU core; we did not use a deep-learning framework. The phase spaces have two and four dimensions, where exact references can be computed by quadrature; in higher dimensions only the learned side and the direct entropy measurement would remain. In four dimensions the nearest-neighbour estimator is biased by a few percent, and only its raw and whitened versions together bracket the prediction. At $\sigma_0 = 0.03$ our networks did not reach the required accuracy at $t = 0$, and we make no claim there; a better allocation of training times near $t = 0$, or a larger network, may be needed. The expansion behind Eq.~\eqref{eq:G} leaves a third-order remainder, visible as a plateau of a few percent for $r \geq \kappa$ at $\sigma_0 = 0.3$.

\emph{Outlook.} The margin is a property of a state seen at a resolution, and the learned score made it measurable from samples alone. Two directions follow. One is molecular: the strain rates of a molecular force field and the noise schedules of generative models for molecular configurations are both defined in phase space, and the band~\eqref{eq:band} sets a lower limit on the noise rate of any such model that keeps the physical drift. The other is a uniform version of the gap, a bound $\kappa - R \geq c\sigma$ for all states and not only Gaussian ones, which the measurements here would allow to test on learned states.
